\subsection{I-31. Full matrix cones give a pair budget controlled by the number of rank levels}
\label{I-31}

\paragraph{Attribution and statement.}
The full-matrix cone method in Liu (2026), Section~3.3,
printed pp.~7--8, avoids treating compound matrices as
arbitrary vectors.  We prove an original-coordinate version
with a somewhat looser block count, so no first-block
function-field factorization is needed.  This is a further
checked specialization of that explicitly unrefereed source,
not a claim of historical novelty or a certification of its
affine-orbit theorem.

Fix an allowed rank set
$\mathcal R=\{s_1<\cdots<s_L=w\}\subseteq\{1,\ldots,w\}$,
where $w\ge2$.  Put
\[
 d_w=w^2-1,\qquad
 d_s=(2ws-1)\left(2\max_{t\in\mathcal R,\ t>s}d_t+1\right)
       \quad(s\in\mathcal R,\ s<w),
 \qquad K_{\mathcal R}=\sum_{s\in\mathcal R}d_s.
 \tag{I.50}
\]
For constant-boundary width-$w$ ROABPs in one hidden order
whose nonempty internal intervals have generic ranks in
$\mathcal R$, the map $t\mathbf1+\operatorname{SV}_{K_{\mathcal R}}$
is a hitting generator in every characteristic.  It satisfies
\[
 K_{\mathcal R}\le2w^2(4w^2)^{L-1}.
 \tag{I.51}
\]
Thus a fixed number of allowed rank levels gives a polynomial
pair budget in $w$, even with arbitrarily many noncommuting
singular factors.  The sample count is still
\[
 (D+1)\bigl((n-1)D+1\bigr)^{K_{\mathcal R}}
                  \binom{D+K_{\mathcal R}}{K_{\mathcal R}},
 \tag{I.52}
\]
for degree bound $D$, with the working-field convention of
I-21.  For the full rank set $\{1,\ldots,w\}$ this is only
$K_{\mathcal R}=2^{O(w\log w)}$, so it proves polynomial size
for each fixed width, not polynomial dependence on variable
width.  I-29 and I-30 give substantially smaller bounds in
their respective special cases.

\paragraph{1. Cones and composition.}
Work over an algebraic closure of the coefficient field.
For a polynomial matrix $M$, write $\mathcal C(M)$ for the
Zariski closure of all scalar multiples of its values.
This is an irreducible affine cone.  Equality of such cones
preserves every homogeneous polynomial identity in the matrix
entries.  It also composes under multiplication on disjoint
variable sets: replacing a factor by a restriction with the
same cone leaves the cone of the whole product unchanged.
Indeed, taking a dense set in either input of the regular
matrix-multiplication map does not change the closure of its
image, and scalar multiples are absorbed by the cone.
These observations remain true after an extension of the
coefficient field.

\paragraph{2. Deleting stable rank-$s$ blocks.}
Let $B_1,\ldots,B_m$ be polynomial $w\times w$ matrices in
disjoint variable blocks, each of generic rank $s$ and with
$B_i(0)=P_s=\operatorname{diag}(I_s,0)$.
For a constant $a\times w$ matrix $R$ with $\pi R\ne0$,
where $\pi$ selects the first $s$ columns, put
\[
 M_i=RB_1\cdots B_i,\qquad
 V_i=\mathcal C(\pi M_i),\qquad V_0=\mathcal C(\pi R).
\]
Then $V_{i-1}\subseteq V_i$, since the zero value of $B_i$
preserves the first $s$ columns.  The cones are irreducible,
start at a line, and lie in the $as$-dimensional matrix space.
There are at most $as-1$ strict dimension increases.

Suppose $V_{j-1}=V_j=V_{j+1}=V$.  We claim that setting
block $j$ to zero preserves the full cone of $M_{j+1}$.
Fix a value of $B_{j+1}$ on the dense open where its top-left
$s\times s$ block $E$ is invertible, and call its top-right
block $F$.  Rank $s$ gives the row-graph identity: any matrix
obtained by multiplying $B_{j+1}$ on the left is determined
by its first $s$ columns $Q$, and equals $Q[I_s,E^{-1}F]$.
After the restriction $B_j=P_s$, those first columns have
cone $VE$.  This is contained in $V$, since the restricted
values are among the original values contributing to
$V_{j+1}=V$.  But $VE$ is closed and irreducible, has
dimension $\dim V$, and is contained in the irreducible $V$;
therefore $VE=V$.

Every original output at this fixed value of $B_{j+1}$ has
its first columns in $V$, and hence belongs to the restricted
full cone by the row-graph identity.  Taking the closure over
the dense set of such values of $B_{j+1}$ proves the claim.
The opposite containment is automatic from restriction.
It also preserves all subsequent cones by composition.

Retain every strict-growth block and the last block of each
run of nongrowth blocks.  All other blocks can be deleted
one at a time using the claim, while the projected cones
and the full cone at the end of each affected pair remain
unchanged.  The number retained is at most
$2(as-1)+1=2as-1$.  In particular, with $R=I_w$, at most
$2ws-1$ blocks preserve the full product cone.

\paragraph{3. Coordinate compression of a preserved-rank word.}
We prove that if every interval of a univariate matrix word
has its generic rank at zero, all these ranks belong to
$\mathcal R$, and the full word has rank $s$, then setting
all but at most $d_s$ positions to zero preserves its full
matrix cone.

For $s=w$, all constant factors are invertible.  Gauge them
to identities.  The prefix cones are nested, starting with
the line through $I_w$ in dimension $w^2$.  Retain only the
strict-growth positions.  At a nongrowth position, replacing
that factor by its identity constant preserves the prefix
cone and then all later cones by composition.  At most
$w^2-1=d_w$ positions remain.

For $s<w$, greedily split the word at successive shortest
prefixes of rank $s$.  Write such a block as $C_jA_{e_j}$;
the preceding chunk $C_j$ is empty or has rank greater
than $s$, as do all of its subintervals.  When the remaining
tail no longer has rank $s$, absorb it into the last block;
that tail is empty or has rank greater than $s$ as well.
Every resulting block has generic and constant rank $s$.
Every product of consecutive blocks has constant rank $s$,
because it is an original interval and cannot have rank
below that of the full word.

Compatible constant bases at the block cuts normalize all
block constants to $P_s$.  To see this, transport an active
$s$-space injectively along the constant products.  At an
internal cut the left kernel of the next block complements
this active space, since the adjacent product retains rank
$s$.  Choosing those passive kernel bases makes each block
constant $\operatorname{diag}(I_s,0)$.
These changes multiply the entire word by fixed invertible
matrices at its ends and do not affect coordinate supports
or cone-preserving restrictions.

Part~2 retains at most $2ws-1$ blocks.  In each retained
block, the induction hypothesis compresses the higher-rank
prefix chunk $C_j$ to at most
$D_s=\max_{t>s,\ t\in\mathcal R}d_t$ positions; the final
tail, if present, costs at most another $D_s$.  Retain the
one factor $A_{e_j}$ between them.  By disjoint-block
composition these restrictions preserve the block cone and
then the full word cone.  Thus at most
$(2ws-1)(2D_s+1)=d_s$ positions suffice, completing the
induction.

\paragraph{4. Absorbing rectangular endpoints in the ground space.}
For an $r\times w$ polynomial matrix $L(Y)$ whose constant
$L_0$ has full row rank $r$, choose a constant right inverse
$R_0$ with $L_0R_0=I_r$.  Then
\[
 B(Y)=I_w+R_0(L(Y)-L_0),\qquad
 L(Y)=L_0B(Y),\qquad B(0)=I_w.
 \tag{I.53}
\]
The determinant of $B$ has constant term one, so $B$ is
generically invertible.  Similarly a $w\times r$ endpoint
$V(Z)$ with full-column-rank constant $V_0$ can be written
$V(Z)=C(Z)V_0$, where $C(0)=I_w$ and $C$ is generically
invertible, by choosing a constant left inverse of $V_0$.
Thus these endpoint factors can be absorbed in the original
$w$-dimensional matrix word, without passing to a compound
space.  Generic and constant interval ranks are unchanged,
just as in I-30.

\paragraph{5. Rank descent and final scalar detection.}
Start with the common center $X_i=t$, which preserves every
full-rank interval.  Process the remaining ranks in
$\mathcal R$ in decreasing order.  For a rank-$r$ local test
as in I-21, the interior chunk has rank greater than $r$,
and every one of its intervals is already preserved.
A selected nonzero $r$-minor of the test has the form
\[
 \det\bigl(L(Y_0)M(\mathbf Y)V(Y_{m+1})\bigr).
\]
The selected endpoint matrices have full row and column
rank $r$ at the current center.  Indeed, a nonzero selected
minor implies that each endpoint has some nonzero
univariate $r$-minor, and each current coordinate remains
individually transcendental over the original coefficient
field: it is $t$ plus terms independent of $t$.

Absorb these endpoints by (I.53).  The enlarged interior
word still has precisely the preserved higher-rank profile.
If its full rank is $s>r$, Part~3 gives a restriction to at
most $d_s\le D_r$ positions preserving its full matrix cone.
The fixed-endpoint determinant is homogeneous of degree
$r$ in the full matrix entries, so this restriction keeps
the selected minor nonzero.  It therefore has a monomial
of support at most $D_r$ and is hit by that many fresh SV
pairs.  All chosen test minors remain simultaneously nonzero
because their product is a nonzero polynomial.  Earlier
ranks remain preserved by setting the new amplitudes to zero.
The local-rank induction proves preservation of all internal
intervals at this rank.

After the smallest allowed rank is processed, every internal
interval is preserved.  Part~3 applied to the full word gives
a cone-preserving restriction on at most $d_s$ positions,
where $s$ is its generic rank.  The nonzero scalar output
is a homogeneous linear functional on the matrix entries,
so it remains nonzero on that restriction.  At most
$\max_{s\in\mathcal R}d_s=d_{s_1}$ additional SV pairs hit it.

The numbers $d_s$ strictly increase as the allowed rank
decreases.  At rank $s_i$ the preservation cost is therefore
$d_{s_{i+1}}$.  Summing those costs and the last scalar
detection cost gives exactly $K_{\mathcal R}$ in (I.50).
All pairs are fresh and their additive maps combine into
one SV map.  The map depends only on the bounds and the
chosen nodes, not on the unknown order or matrices.

Finally $d_s+1\le4w^2(D_s+1)$ and $d_w+1=w^2$.
The recurrence also gives $d_s>2D_s$ for $s<w$.
Hence the sum is at most twice its largest term, yielding
(I.51).  The degree and Newton-simplex sampling argument is
identical to I-21, proving (I.52).
\hfill$\square$

\paragraph{Remaining main gate.}
The proof replaces the exponential compound dimension by
matrix-cone dimensions, but a new block-selection factor
appears at each distinct singular rank level.  For bounded
$L$ this is polynomial in $w$ at the pair-count level;
for unrestricted $L$ the bound is still exponential in
$w\log w$.  Moreover even a polynomial pair budget would
still need an efficient sampling argument to yield a
polynomial-size hitting set.  No implication from Hoza's
path isolation to either of these missing steps has been
proved.  The dense lower bound I-26, the filtered-mask
obstruction I-27, and the known logarithmic every-order
result I-B08 continue to delimit the candidate repairs.
